\uuid{Gn6E}
\titre{Étude de fonction}
\auteur{bodin ; hanani ; mzari}
\organisation{exo7}
\licence{CC BY-NC-SA 4.0}
\datecreate{2020-01-20}
\niveau{L1}
\module{Analyse}
\chapitre{Fonctions circulaires et hyperboliques inverses}
\sousChapitre{Autre}
\difficulte{3}
\theme{Fonctions usuelles, étude de fonction}
\exo7qcmid{lille-1.fonctions-usuelles.021}

\contenu{
\begin{question}
\qclassification{126.00}
\qtags{motcle=étude de fonction}
Soit $(E)$ l'équation : $ x^x=(\sqrt x)^{x+1}$.  Quelles sont les assertions vraies ?

\begin{answers}

     \good{Le domaine de définition de $(E)$ est $]0,+\infty[$.}

    \bad{$(E)$ n'admet pas de solution.}

    \bad{$(E)$ admet deux solutions distinctes.}

    \good{$(E)$ admet une unique solution.}
\end{answers}
\begin{explanations}
$(E)$ est définie si $x>0$.

Soit $x>0$, alors $(E) \Leftrightarrow x \ln x  = \frac{1}{2}(x+1)\ln x  \Leftrightarrow (x-1) \ln x = 0 \Leftrightarrow x=1.$
\end{explanations}
\end{question}
}
